\documentclass[11pt]{amsart}
\usepackage[margin=1in]{geometry}
\usepackage{amsmath,amssymb,mathtools}
\usepackage{graphicx}
\usepackage{xcolor}
\usepackage{hyperref}
\hypersetup{colorlinks=true,linkcolor=blue!55!black,citecolor=blue!55!black,urlcolor=blue!55!black}
\emergencystretch=2em

\newtheorem{theorem}{Theorem}
\newtheorem{lemma}[theorem]{Lemma}
\newtheorem{corollary}[theorem]{Corollary}
\theoremstyle{definition}
\newtheorem{definition}[theorem]{Definition}
\newtheorem{remark}[theorem]{Remark}

\newcommand{\Dsp}{\mathcal D}
\newcommand{\Ber}{\operatorname{Ber}}
\newcommand{\ber}{\operatorname{ber}}
\newcommand{\Num}{W}
\newcommand{\Disk}{\overline{\mathbb D}}

\title[Dirichlet Berezin-sectorial shifts]{Explicit Dirichlet composition--differentiation shifts\\
resolve both Berezin-sectorial construction alternatives}
\author{Candidate solution packet for arXiv:2601.02817}
\date{13 August 2026}

\begin{document}
\begin{abstract}
The conclusion of Mahapatra--Mukherjee--Sen--Birbonshi--Paul asks whether
scalar shifts of composition--differentiation operators on the Dirichlet
space can systematically produce (i) Berezin sectorial but nonsectorial
operators or (ii) operators that are sectorial in both senses with different
indices.  We give both constructions from one symbol.  For
\(\phi(z)=z/2\) and \(Df=f'\circ\phi\), the Berezin range is a disk of radius
\(b<9/20\), whereas the numerical range contains the disk of radius
\(\sqrt{7/24}>9/20\).  Thus \(aI+D\) is Berezin sectorial but not sectorial
whenever \(9/20<a<\sqrt{7/24}\).  Moreover, \(I+D\) is sectorial and its
classical index is strictly larger than its Berezin index.  The only finite
calculation in the proof is supplied as an exact rational Bernstein
certificate.
\end{abstract}

\maketitle

\noindent
\fbox{\parbox{0.94\textwidth}{\textbf{Status.}
Candidate full answer, likely valid, subject to human review.  The result
answers both alternatives of the intended \emph{shifted}
composition--differentiation construction problem.}}

\section{Source problem and scope}

The source is S.~Mahapatra, S.~Mukherjee, A.~Sen, R.~Birbonshi, and K.~Paul,
\emph{An introduction of Berezin sectorial operators and its application to
Berezin number inequalities}, arXiv:2601.02817 (2026)~\cite{source}.
On PDF page 22 the authors explicitly motivate the question by ``applying
certain shifts to composition-differentiation operators.''  The conclusion on
PDF page 23 asks whether such operators can be systematically constructed
``to be Berezin sectorial but not sectorial or both Berezin sectorial and
sectorial but with different indices.''

\begin{center}
\includegraphics[width=0.98\textwidth]{figures/open_problem_crop.png}

\small Source evidence: conclusion, PDF page 23 of arXiv:2601.02817.
\end{center}

We answer both alternatives with scalar shifts of a single
composition--differentiation operator, precisely in the sense of the preceding
source paragraph.  We do not claim that \(aI+D_\phi\) itself can be rewritten
as one unshifted \(D_\psi\).

\section{Definitions}

We use the Dirichlet Hilbert space and norm fixed by the source:
\[
 \Dsp=\left\{f(z)=\sum_{n\ge0}a_nz^n:
        \|f\|_{\Dsp}^{2}=\sum_{n\ge0}(n+1)|a_n|^2<\infty\right\}.
\]
Thus
\[
 e_n(z)=\frac{z^n}{\sqrt{n+1}}\quad(n\ge0)
\]
is an orthonormal basis.  The reproducing kernel at \(\lambda\in\mathbb D\)
and its squared norm are
\[
 k_\lambda(z)=\sum_{n\ge0}\frac{\overline\lambda^n z^n}{n+1},
 \qquad
 K(t):=\|k_\lambda\|^2=\sum_{n\ge0}\frac{t^n}{n+1}
       =\frac{-\log(1-t)}{t},
 \qquad t=|\lambda|^2,
\]
with the value at \(t=0\) understood by continuity.

For a bounded operator \(T\) on \(\Dsp\),
\[
 \Num(T)=\{\langle Tx,x\rangle:\|x\|=1\},\qquad
 \Ber(T)=\{\langle T\widehat k_\lambda,\widehat k_\lambda\rangle:
                 \lambda\in\mathbb D\}.
\]
For \(0\le\theta<\pi/2\), put
\(S_\theta=\{z:|\arg z|\le\theta\}\).  The operator is sectorial if
\(\Num(T)\subseteq S_\theta\), and Berezin sectorial if
\(\Ber(T)\subseteq S_\theta\), for some such \(\theta\).  We denote the
infima of admissible semi-angles by \(\theta_{\rm num}(T)\) and
\(\theta_{\rm Ber}(T)\), respectively.

For the analytic self-map \(\phi(z)=z/2\), let
\[
 D=D_\phi,\qquad (Df)(z)=f'(\phi(z))=f'(z/2).
\]

\section{Idea of the proof}

In the monomial basis, \(D\) is a compact backward weighted shift.  This has
two complementary consequences.  First, its Berezin symbol can be summed
radially: it is \(\overline\lambda\) times a positive function of
\(|\lambda|^2\), so its Berezin range is exactly a centered disk.  An exact
polynomial majorant, certified in Bernstein form, puts that disk inside radius
\(9/20\).  Second, the first two shift weights already produce, on a
three-dimensional compression, a numerical circle of radius
\(\sqrt{7/24}\).  The rational threshold \(9/20\) lies strictly between the
two radii.  Translating both ranges by \(a>0\) then yields the desired
sectorial separations.  A uniform norm bound below one makes the shift by
\(I\) classically sectorial as well, but with a strictly larger index.

\section{The construction theorem}

Define
\[
 F(t)=\sum_{n\ge1}\frac{n}{n+1}\left(\frac t2\right)^{n-1},
 \qquad
 b=\max_{0\le t\le1}\frac{\sqrt t\,F(t)}{K(t)},
\]
where the quotient is extended continuously by zero at \(t=1\).

\begin{theorem}\label{thm:main}
For \(D f=f'(z/2)\) on \(\Dsp\), the following hold.
\begin{enumerate}
\item \(\Ber(D)=\{z:|z|\le b\}\) and \(b<9/20\).
\item \(\Num(D)\) contains \(\{z:|z|\le\sqrt{7/24}\}\).
\item For every
\[
 \frac9{20}<a<\sqrt{\frac7{24}},
\]
the shifted composition--differentiation operator \(T_a=aI+D\) is Berezin
sectorial but not sectorial.
\item The operator \(T_1=I+D\) is both sectorial and Berezin sectorial, and
its optimal indices satisfy
\[
 \theta_{\rm Ber}(T_1)=\arcsin b
 <\arcsin\frac9{20}
 <\arcsin\sqrt{\frac7{24}}
 \le \theta_{\rm num}(T_1)<\frac\pi2.
\]
\end{enumerate}
Consequently one explicit symbol gives systematic families resolving both
alternatives in the source problem.
\end{theorem}

\begin{proof}
Acting on the orthonormal basis gives
\begin{equation}\label{eq:weights}
 De_0=0,\qquad
 De_n=\alpha_ne_{n-1},\qquad
 \alpha_n=n\sqrt{\frac n{n+1}}\,2^{1-n}\quad(n\ge1).
\end{equation}
In particular \(\alpha_n\to0\), so \(D\) is compact and bounded.

Using the convention that the inner product is linear in its first variable,
the kernel expansion and~\eqref{eq:weights} yield
\[
 \langle Dk_\lambda,k_\lambda\rangle
 =\overline\lambda\sum_{n\ge1}\frac n{n+1}
       \left(\frac{|\lambda|^2}{2}\right)^{n-1}
 =\overline\lambda F(t).
\]
Hence
\begin{equation}\label{eq:berezin-symbol}
 \widetilde D(\lambda)
 =\overline\lambda\frac{F(t)}{K(t)},\qquad t=|\lambda|^2.
\end{equation}
The radial modulus in~\eqref{eq:berezin-symbol} is continuous, vanishes at
\(t=0\), and also tends to zero as \(t\to1^-\), because \(F(1)<\infty\)
while \(K(t)\to\infty\).  It therefore attains a maximum \(b\).  Every
angle occurs, and the intermediate value theorem supplies every radius
between zero and \(b\).  Thus \(\Ber(D)=b\Disk\).

It remains to prove the strict rational bound for \(b\).  For \(0\le t\le1\),
set
\begin{align*}
 P(t)&=\sum_{n=0}^{8}\frac{t^n}{n+1},\\
 N(t)&=\sum_{n=1}^{8}\frac n{n+1}\left(\frac t2\right)^{n-1}
        +2\left(\frac t2\right)^8.
\end{align*}
Plainly \(K(t)\ge P(t)\).  Also
\[
 \sum_{n\ge9}\frac n{n+1}\left(\frac t2\right)^{n-1}
 \le\frac{(t/2)^8}{1-t/2}\le2(t/2)^8,
\]
so \(F(t)\le N(t)\).  Consider the exact degree-17 polynomial
\begin{equation}\label{eq:q}
 q(t)=\frac{81}{400}P(t)^2-tN(t)^2.
\end{equation}
Write \(q(t)=\sum_{j=0}^{17}q_jt^j\).  Its degree-32 Bernstein
coefficients on \([0,1]\) are
\[
 \beta_k=\sum_{j=0}^{\min(k,17)}q_j
          \frac{\binom{k}{j}}{\binom{32}{j}},\qquad 0\le k\le32.
\]
Exact rational expansion gives \(\beta_k>0\) for all 33 indices, with
\[
 \min_{0\le k\le32}\beta_k
 =\beta_{28}
 =\frac{55075026511}{201587097600000}>0.
\]
Because the Bernstein basis functions are nonnegative and sum to one,
\(q(t)>0\) throughout \([0,1]\).  The attached checker performs this finite
expansion using rational arithmetic only.  From~\eqref{eq:q},
\[
 \frac{\sqrt t\,F(t)}{K(t)}
 \le\frac{\sqrt t\,N(t)}{P(t)}<\frac9{20},
\]
and therefore \(b<9/20\).

For the numerical lower bound, compress \(D\) to
\(E=\operatorname{span}\{e_0,e_1,e_2\}\).  The first two weights are
\[
 \alpha_1=\frac1{\sqrt2},\qquad
 \alpha_2=\sqrt{\frac23}.
\]
The largest eigenvalue of the real part of this three-by-three weighted shift
is
\[
 r_0=\frac12\sqrt{\alpha_1^2+\alpha_2^2}
     =\frac12\sqrt{\frac12+\frac23}
     =\sqrt{\frac7{24}}.
\]
Thus \(r_0\in\Num(D)\).  If \(U_\vartheta e_n=e^{in\vartheta}e_n\), then
\(U_\vartheta^*DU_\vartheta=e^{i\vartheta}D\).  Unitary invariance gives
\(r_0e^{i\vartheta}\in\Num(D)\) for every \(\vartheta\).  Convexity of the
numerical range (Toeplitz--Hausdorff; see~\cite{gustafsonrao}) now gives
\(r_0\Disk\subseteq\Num(D)\).  The two radii are strictly separated because
\[
 \frac7{24}-\left(\frac9{20}\right)^2=\frac{107}{1200}>0.
\]

Translations satisfy
\[
 \Ber(aI+D)=a+\Ber(D),\qquad \Num(aI+D)=a+\Num(D).
\]
For \(9/20<a<r_0\), the first set is the disk of center \(a\) and radius
\(b<a\), hence lies in the sector of semi-angle \(\arcsin(b/a)<\pi/2\).
But \(a-r_0<0\) belongs to the second set, so no sector \(S_\theta\) with
\(\theta<\pi/2\) can contain it.  This proves assertion~3.

Finally, the norm of a unilateral weighted shift is the supremum of its
weights.  Here \(\alpha_2=\sqrt{2/3}\), while \(\alpha_1<\sqrt{2/3}\) and,
for \(n\ge3\),
\[
 \alpha_n\le\frac{n}{2^{n-1}}\le\frac34<\sqrt{\frac23}.
\]
Consequently \(\|D\|=\sqrt{2/3}<1\).  Therefore
\[
 \Num(I+D)\subseteq\{z:|z-1|\le\sqrt{2/3}\},
\]
a disk contained in a sector of semi-angle strictly below \(\pi/2\); hence
\(I+D\) is sectorial.  Since \(\Ber(I+D)=1+b\Disk\), elementary tangent
geometry gives \(\theta_{\rm Ber}(I+D)=\arcsin b\).  Since
\(1+r_0\mathbb T\subseteq\Num(I+D)\), the same tangent geometry gives
\(\theta_{\rm num}(I+D)\ge\arcsin r_0\).  The strict chain in assertion~4
follows.
\end{proof}

\section{Exact verification}

Run from the packet directory:
\begin{verbatim}
conda run --no-capture-output -n sandbox python \
  code/verify_radius_gap.py
\end{verbatim}
The script constructs \(P,N,q\) from their displayed definitions, converts
\(q\) to degree-32 Bernstein form, checks all 33 coefficients with
\texttt{fractions.Fraction}, confirms the exact minimum displayed above, and
checks the exact gap \(107/1200\).  Its terminal line is
\begin{verbatim}
CERTIFIED: b < 9/20 < sqrt(7/24)
\end{verbatim}
This is an exact finite certificate, not a floating-point sampling argument.
An exploratory decimal estimate \(b\approx0.433904\) was used only to select
the convenient rational separator and is not used in the proof.

\section{Novelty check, limitations, and review recommendation}

On 12 August 2026 we searched the run's registry, solution, attempt, and
proof-gap indexes for arXiv:2601.02817 and the phrases \emph{Berezin
sectorial}, \emph{composition-differentiation}, and \emph{Dirichlet space};
no duplicate or earlier result was found.  A bounded arXiv web search used the
source id and the combinations ``Berezin sectorial'' with
``composition-differentiation'', ``Dirichlet'', and ``different indices''.
It found the source paper but no later paper answering this construction
question.  This is a bounded search rather than an exhaustive MathSciNet,
zbMATH, or citation-graph review, so novelty confidence is moderate rather
than definitive.

The result is full for the source's either/or construction request and in fact
supplies both alternatives.  It does not classify all symbols, compute the
exact constants \(b\) or \(w(D)\), or address unshifted \(D_\psi\) without
the scalar translations explicitly contemplated by the source.  Human review
should focus on (i) the Dirichlet normalization, (ii) the Berezin-symbol
calculation, (iii) the exact Bernstein conversion, and (iv) agreement that the
source problem concerns scalar shifts.  Subject to those checks, promotion as
a full affirmative answer is recommended.

\begin{thebibliography}{9}
\bibitem{source}
S.~Mahapatra, S.~Mukherjee, A.~Sen, R.~Birbonshi, and K.~Paul,
\emph{An introduction of Berezin sectorial operators and its application to
Berezin number inequalities}, arXiv:2601.02817, 2026.

\bibitem{gustafsonrao}
K.~E.~Gustafson and D.~K.~M.~Rao,
\emph{Numerical Range: The Field of Values of Linear Operators and Matrices},
Springer, New York, 1997.
\end{thebibliography}

\end{document}
